\chapter{從普通對稱性到格拉斯曼代數}

\lectureoverview{超對稱是一種對稱性，因此我們先回顧這個詞在量子力學中的意義。物理變換由作用於希爾伯特空間的么正算符表示；連續變換可由無窮小變換生成，而其交換子則編碼了變換的組合方式。真正的對稱性會保持能量不變，這等價於變換與哈密頓量對易。超對稱改變了所有熟悉例子共有的一項假設：它的生成元會將玻色子變為費米子，也會將費米子變為玻色子。構造這類生成元，自然會引出產生與湮滅算符、反對易子，以及奇元素平方為零的有限維格拉斯曼代數。}

\section{對稱性關聯不同的組態}

對稱性不只是裝飾性的規律，而是一種將一個物理組態與另一個組態關聯起來的操作。例如，左右反射會將兩隻手彼此對應。如果一個人及其周遭定律具有精確的左右對稱性，兩隻手的相應性質就會一致。更一般地說，由精確對稱性相連的物體或態，必須具有該對稱性所保持的物理量。在粒子物理學中，這通常意味著質量相等；在原子物理學中，則意味著同一轉動多重態中的態具有相同能量。

我們的目標是超對稱，但不妨先回顧轉動和平移等連續對稱性。超對稱在通常意義上是否為「連續」對稱性，這個問題本身就不容易回答。它的無窮小參數不是普通的可交換實數。最好將它視為連續對稱性的分次推廣；下文引入的格拉斯曼算術，足以讓這個說法在目前的討論中變得精確。\editorialnote{現代術語將超對稱描述為連續分次對稱性。這項限定說明了講座中的提醒：通常的連續與離散之分，並不完全適用於超對稱。}

\section{希爾伯特空間上的么正操作}

量子系統以態向量\(\ket{\psi}\)表示。物理對稱操作則以作用於該態的算符\(U\)表示：
\begin{equation}
\ket{\psi}\longmapsto U\ket{\psi}.
\end{equation}
對空間轉動而言，轉軸與轉角決定了\(U\)。如果一個態包含位於多個位置的粒子，變換後的右矢便代表這些粒子都位於各自旋轉後的位置。平移的作用方式類似：它將原態變為整個組態都已平移的態。

\begin{classroomqa}
\audiencequestion{態向量屬於一個往往是無限維的希爾伯特空間。我們轉動的是希爾伯特空間，還是普通的三維空間？\lecturetimestamp{00:04:20}}
\lecturerresponse{這兩種說法是在不同層次上描述同一個操作。么正變換作用於希爾伯特空間中的向量，表示組態在普通空間中的物理轉動。這不是字面上將抽象的希爾伯特空間作三維轉動，而是在希爾伯特空間中表示空間轉動。}
\end{classroomqa}

關鍵條件是么正性：
\begin{equation}
U^\dagger U=UU^\dagger=\mathbf 1,
\qquad
U^{-1}=U^\dagger.
\end{equation}
這會保持內積，因而保持機率：
\begin{equation}
\braket{U\psi}{U\psi}
=\bra{\psi}U^\dagger U\ket{\psi}
=\braket{\psi}{\psi}.
\end{equation}

\begin{figure}[H]
\centering
\includegraphics[width=0.56\textwidth]{lecture_04_figure_02.png}
\lectureframe{00:05:56.420}
\caption{黑板上先從變換後的態過渡到么正性條件。右矢與帶匕首的乘積清晰可見，但講者遮住了等式右側的部分；完整方程見下文推導。}
\end{figure}

嚴格來說，維格納定理也容許反么正對稱性，例如時間反演。本文所討論、與恆等變換連續相連的變換，則由么正算符表示。\editorialnote{反么正變換是標準的補充說明，不影響講座對轉動、平移或超對稱生成元的討論。}

\section{無窮小變換與生成元}

什麼都不做所對應的是\(U=\mathbf 1\)，而不是\(U=0\)。足夠小的變換與恆等變換只差一個對小參數呈線性的項。對於轉過小角度\(\epsilon\)的轉動，
\begin{equation}
U(\epsilon)=\mathbf 1+i\epsilon L+O(\epsilon^2).
\end{equation}
么正性在一階近似下要求
\begin{align}
U^\dagger U
&=\left(\mathbf 1-i\epsilon L^\dagger\right)
  \left(\mathbf 1+i\epsilon L\right)+O(\epsilon^2)\\
&=\mathbf 1+i\epsilon(L-L^\dagger)+O(\epsilon^2),
\end{align}
因此
\begin{equation}
L^\dagger=L.
\end{equation}
算符\(L\)就是無窮小生成元。

轉動需要三個生成元，分別對應三個獨立的軸：
\begin{equation}
L_x,\qquad L_y,\qquad L_z.
\end{equation}
繞任意軸\(\widehat{\bm n}\)的轉動由下式生成：
\begin{equation}
L_{\widehat{\bm n}}=\widehat{\bm n}\cdot\bm L.
\end{equation}

\begin{classroomqa}
\audiencequestion{繞\(x\)、\(y\)和\(z\)軸的生成元，可以相加來生成繞另一軸的轉動嗎？\lecturetimestamp{00:07:50}}
\lecturerresponse{可以。它們的線性組合能生成繞中間方向軸的轉動。三個獨立分量足以構成所有無窮小空間轉動。}
\end{classroomqa}

對轉動而言，這些生成元就是角動量的分量。這套記號也適用於更一般的情況，因為任何連續對稱性都可以在局部透過其生成元來理解。

\section{作為乘法表的交換子}

轉動生成元滿足
\begin{align}
[L_x,L_y]&=iL_z,&
[L_y,L_z]&=iL_x,&
[L_z,L_x]&=iL_y,
\end{align}
簡潔地寫成
\begin{equation}
[L_i,L_j]=i\epsilon_{ijk}L_k.
\end{equation}
兩個指標相同時，Levi-Civita 符號為零；對\(1,2,3\)的偶排列或奇排列，它分別等於\(+1\)或\(-1\)。關鍵的結構性質是封閉性：兩個生成元的交換子仍是生成元。

為什麼交換子能描述轉動的幾何？先繞\(x\)軸作一次小轉動，再繞\(y\)軸作一次小轉動，接著反向轉回\(x\)軸的轉動，最後反向轉回\(y\)軸的轉動。這些操作不會彼此抵消，因為繞不同軸的轉動不對易。

\begin{classroomqa}
\audiencequestion{這種不對易現象，可以用飛機的滾轉、俯仰和偏航來理解嗎？\lecturetimestamp{00:12:36}}
\lecturerresponse{可以。依序繞機體軸轉動，便能直接從幾何上看出改變操作順序會改變最終姿態。下文的代數正是這個熟悉事實的無窮小版本。}
\end{classroomqa}

令
\begin{equation}
U_x(\epsilon)=\mathbf 1+i\epsilon L_x,
\qquad
U_y(\delta)=\mathbf 1+i\delta L_y.
\end{equation}
由於算符從右向左作用於右矢，所述的操作順序可表示為
\begin{equation}
\mathcal C
=U_y^{-1}(\delta)U_x^{-1}(\epsilon)
 U_y(\delta)U_x(\epsilon).
\end{equation}

\begin{classroomqa}
\audiencequestion{為什麼反向操作可以寫成逆算符，而不是厄米共軛？\lecturetimestamp{00:13:25}}
\lecturerresponse{因為這些變換是么正的：\(U^\dagger=U^{-1}\)。對無窮小變換取逆，只需反轉一階項的符號：\(U_x^{-1}=\mathbf 1-i\epsilon L_x+O(\epsilon^2)\)。}
\end{classroomqa}

在黑板上計算這四個因子的乘積，成了全班一起投入的練習：所有關於\(\epsilon\)或\(\delta\)的一階項都抵消了，大家逐一追蹤混合項；一度出現的係數二混淆，也顯示算符順序很容易弄錯。整理後的二階結果是
\begin{align}
\mathcal C
&=\mathbf 1+\epsilon\delta[L_x,L_y]
  +O(\epsilon^2\delta,\epsilon\delta^2)\\
&=\mathbf 1+i\epsilon\delta L_z
  +O(\epsilon^2\delta,\epsilon\delta^2).
\end{align}
因此，這個閉合操作序列留下了一個繞\(z\)軸的小轉動，其轉角正比於\(\epsilon\delta\)。\editorialnote{經過課堂上的嘗試性計算後，這裡列出的乘法結果確定了操作順序與係數。若採用相反的序列約定，符號也會反轉；但殘餘變換由交換子決定這一結論不變。}

\begin{classroomqa}
\audiencequestion{如果改變反向操作的順序，乘積會變成恆等變換嗎？\lecturetimestamp{00:25:23}}
\lecturerresponse{如果每個變換都緊接在它自己的逆變換旁邊，成對的因子就會抵消。這個非平凡的交換子序列刻意將它們分開；它無法化為恆等變換，正是順序依賴性的表現。}
\end{classroomqa}

較長的乘積確實會更複雜，但沒有引入新的原理。

\begin{classroomqa}
\audiencequestion{如果生成元多於三個，或有許多小變換，會發生什麼事？\lecturetimestamp{00:25:34}}
\lecturerresponse{交換子會更多，計算也更繁複；但只要知道封閉代數，就能確定任意多個無窮小變換如何組合。有限變換可以由許多小變換組成。}
\end{classroomqa}

就此而言，交換子代數就是連續群在恆等元附近的乘法表。

\section{平移補足空間變換的例子}

空間平移由動量生成。依照講座採用的符號約定，
\begin{equation}
\begin{aligned}
T_x(a)&=\mathbf 1+iaP_x+O(a^2),\\
T_y(b)&=\mathbf 1+ibP_y+O(b^2),
\end{aligned}
\end{equation}
其中\(a\)和\(b\)代表距離，而不是角度。

兩個平移彼此對易。先將一點沿\(x\)方向移動，再沿\(y\)方向移動；接著反向移動\(x\)方向的平移，最後反向移動\(y\)方向的平移，該點便回到原點：
\begin{equation}
[P_i,P_j]=0.
\end{equation}

\begin{classroomqa}
\audiencequestion{平移或轉動物理物體，是否需要施力才能讓它開始和停止運動？\lecturetimestamp{00:28:45}}
\lecturerresponse{這裡所定義的操作並非如此。對稱變換是比較原始狀態與整個組態已平移或轉動後的數學狀態。它描述的不是動力學軌跡、加速度，也不是移動裝置時所做的功。}
\end{classroomqa}

轉動和平移不對易。先平移一個點，再令它繞原點轉動，接著沿原來的軸反向平移，最後反向轉回；該點不會回到原來的位置。混合代數關係為
\begin{equation}
[L_i,P_j]=i\epsilon_{ijk}P_k.
\end{equation}
這個關係表示動量本身在轉動下會如向量般變換。此處繁複的指標計算不是重點；一般性的結論是，歐幾里得空間的變換由轉動和平移生成元之間的交換子編碼。

\section{對稱性必須保持能量}

現在我們從變換的組合轉向動力學。如果\(U\)確實是系統的對稱變換，施行它就不能改變能量。設\(\ket{E}\)為能量本徵態：
\begin{equation}
H\ket{E}=E\ket{E}.
\end{equation}
等式右側的\(E\)是數；\(\ket{E}\)則是以該本徵值標記的向量。

\begin{figure}[H]
\centering
\begin{minipage}[t]{0.43\textwidth}
\centering
\includegraphics[width=\linewidth]{lecture_04_figure_03.png}
\lectureframe{00:34:38.239}
\end{minipage}
\hfill
\begin{minipage}[t]{0.50\textwidth}
\centering
\includegraphics[width=\linewidth]{lecture_04_figure_04.png}
\lectureframe{00:36:22.376}
\end{minipage}
\caption{變換後的右矢\(U\ket{E}\)與定義本徵值的方程\(H\ket{E}=E\ket{E}\)，分別呈現在連續的黑板畫面中。}
\end{figure}

如果\(U\)是對稱變換，那麼\(U\ket{E}\)也具有相同的本徵值：
\begin{equation}
HU\ket{E}=EU\ket{E}.
\end{equation}
由於\(E\)是純量，與\(U\)對易，因此可以使用原本的本徵值方程：
\begin{align}
HU\ket{E}
&=EU\ket{E}\\
&=UE\ket{E}\\
&=UH\ket{E}.
\end{align}
所以
\begin{equation}
(HU-UH)\ket{E}=0.
\end{equation}

\begin{classroomqa}
\audiencequestion{能否從第一個關係式重新推導一次？\lecturetimestamp{00:35:55}}
\lecturerresponse{從\(H\ket{E}=E\ket{E}\)開始。對稱性表示\(U\ket{E}\)具有相同能量，因此\(HU\ket{E}=EU\ket{E}\)。將純量\(E\)移過\(U\)，再以\(H\ket{E}\)取代\(E\ket{E}\)，然後相減，便得到\((HU-UH)\ket{E}=0\)。}
\end{classroomqa}

\begin{figure}[H]
\centering
\includegraphics[width=0.70\textwidth]{lecture_04_figure_05.png}
\lectureframe{00:38:51.066}
\caption{完整的黑板推導。彎曲的引導標記將等能量的陳述連到代入式\(E\ket{E}=H\ket{E}\)；最後一行是\((HU-UH)\ket{E}=0\)。}
\end{figure}

這個論證對每一個能量本徵態都成立。假設這些態構成完備基底，那麼對其中每個態作用後都得到零的算符，就是零算符：
\begin{equation}
[H,U]=0.
\end{equation}
這就是與時間無關的么正變換成為哈密頓量對稱變換的判準。

對無窮小變換而言，
\begin{equation}
U(\epsilon)=\mathbf 1+i\epsilon L_i+O(\epsilon^2),
\end{equation}
恆等算符本來就與\(H\)對易，因此每個生成元都必須滿足
\begin{equation}
[L_i,H]=0.
\end{equation}

\begin{figure}[H]
\centering
\includegraphics[width=0.47\textwidth]{lecture_04_figure_06.png}
\lectureframe{00:42:14.296}
\caption{黑板上精簡列出的對稱性判準：\([U,H]=0\)，接著是其生成元形式；後者依據部分模糊的下標整理為\([L_i,H]=0\)。}
\end{figure}

如果\(L_i\)沒有顯含時間，海森堡方程給出
\begin{equation}
\frac{dL_i}{dt}
=i[H,L_i]=0
\qquad(\hbar=1).
\end{equation}
因此，對稱性生成元是守恆荷。空間平移對稱性導出動量守恆；轉動對稱性導出角動量守恆。\editorialnote{顯式時間導數的條件是標準的補充說明；講座陳述守恆律時，已隱含了這項條件。}

\section{變換整個物理系統}

如果只變換裝置的一部分，「轉動是對稱性」這句話就可能遭到誤用。自由自旋可以旋轉而不改變能量。然而，若將它置於固定磁場中，卻只旋轉自旋，其塞曼能量就會改變。這不代表轉動對稱性不成立，而是因為自旋相對於磁場轉動時，磁場仍保持固定。若讓自旋和磁場一起轉動，或是外加磁場本身在該操作下保持不變，能量就不會改變。

氫原子也有同樣的問題。只平移電子而讓質子保持不動，會改變兩者的距離，因而改變庫侖能量。平移對稱性作用於完整的電子—質子系統，讓兩者移動相同的距離。

\begin{classroomqa}
\audiencequestion{外部場本身也必須具有該對稱性嗎？\lecturetimestamp{00:43:09}}
\lecturerresponse{固定的外部場必須已在該操作下保持不變，或者外部場必須和系統其他部分一起變換。自旋相對於未轉動的磁場所作的轉動，形成的是不同的物理組態，並非整個裝置經對稱變換後的組態。}
\end{classroomqa}

原子中的轉動簡併是一個清楚的正面例子。對固定角動量\(l\)而言，轉動會混合磁量子數為
\begin{equation}
m=-l,-l+1,\ldots,l.
\end{equation}
的態。在沒有外加磁場或電場時，多重態中的所有成員能量相同：
\begin{equation}
E_{l,m}=E_l.
\end{equation}
對稱性將不同的態彼此關聯，並使它們具有簡併能階。

\section{超對稱改變了什麼}

迄今討論的各種熟悉對稱性，都保留了玻色子與費米子的區別。轉動會將一個電子變為另一個電子，而不會將它變成光子。色旋轉會混合夸克的三種顏色，但不會將夸克變成膠子。普通的內部對稱性和時空對稱性可以混合不同分量或量子數標籤，但它們的生成元是玻色性的，不會改變費米子宇稱。

超對稱則不同：
\begin{equation}
Q\ket{B}\sim\ket{F},
\qquad
Q\ket{F}\sim\ket{B}.
\end{equation}
它的多重態跨越了玻色子與費米子之間的界線。在精確的超對稱理論中，同一多重態中的夥伴粒子質量相同。這正是前幾講所要求的配對：相互匹配的玻色子圈圖與費米子圈圖，可以抵消標量質量受到的大幅紫外修正。這種結構難以直觀理解，因為沒有普通的空間運動能讓整數自旋態變成半整數自旋態。因此，接下來我們採用代數方法。

\section{超荷的簡單模型}

設\(A,A^\dagger\)為玻色子的湮滅與產生算符，\(C,C^\dagger\)為費米子的湮滅與產生算符。交換兩種粒子的單粒子態，最簡單的厄米算符為
\begin{equation}
Q=C^\dagger A+A^\dagger C.
\end{equation}
第一項湮滅一個玻色子並產生一個費米子；第二項則湮滅一個費米子並產生一個玻色子。若\(Q\)作用於單玻色子態，\(A^\dagger C\)項會消失，因為沒有費米子可供湮滅；此時由\(C^\dagger A\)完成轉換。作用於單費米子態時，兩項的作用則互換。

\begin{classroomqa}
\audiencequestion{這種構造會產生一個同時為玻色子與費米子疊加態的單一粒子嗎？\lecturetimestamp{00:52:37}}
\lecturerresponse{\(Q\)是算符，不是態。它作用於純玻色態時會產生費米態；作用於純費米態時則會產生玻色態。若它線性作用於原本已是一般疊加態的態，便會產生另一個疊加態，就像任何線性算符一樣。}
\end{classroomqa}

接著來檢驗真空態。由於兩個湮滅算符作用於真空態都會得到零，
\begin{equation}
A\ket{0}=0,\qquad C\ket{0}=0,
\end{equation}
因此
\begin{equation}
Q\ket{0}=0.
\end{equation}
結果是零向量，而不是另一個真空態。

這個模型刻意採取示意性的簡化。一個自旋為零的玻色子可以轉換成自旋二分之一費米子的任一自旋分量，因此生成元本身需要帶有旋量指標：
\begin{equation}
Q_\alpha\ket{B}\sim\ket{F_\alpha}.
\end{equation}
在講座簡化的雙分量圖像中，
\begin{equation}
Q_\alpha\ket{F_\beta}\sim\delta_{\alpha\beta}\ket{B}.
\end{equation}

\begin{classroomqa}
\audiencequestion{自旋為零的玻色子只有一個自旋分量，自旋二分之一的費米子卻有兩個；這種情況下，態空間的維度如何能相符？\lecturetimestamp{00:55:02}}
\lecturerresponse{單一不帶指標的\(Q\)並不足夠。還有旋量分量\(Q_\alpha\)，每個分量都能產生一個費米子自旋分量。超荷本身帶有半單位自旋，因為它會使自旋改變半個單位。}
\end{classroomqa}

\(Q_\alpha\)是費米性算符，更精確地說，是格拉斯曼奇算符。每一項都含有奇數個費米子算符，因此會反轉費米子宇稱。一般可觀測量則是偶的。必須區分奇超荷與奇變換參數：有限的超對稱變換由\(\epsilon^\alpha Q_\alpha+\bar\epsilon_{\dot\alpha}\bar Q^{\dot\alpha}\)這類偶組合構成，其中\(\epsilon\)和\(Q\)都是奇的。\editorialnote{講座非正式地將生成元稱為格拉斯曼數。更精確地說，超荷是奇算符，反對易參數則是格拉斯曼數；兩者的乘積為偶，可以出現在有限變換中。}

\section{玻色子與費米子的算符代數}

對以\(i,j\)標記的玻色模式，
\begin{align}
[A_i,A_j]&=0,&
[A_i^\dagger,A_j^\dagger]&=0,&
[A_i,A_j^\dagger]&=\delta_{ij}.
\end{align}
這些標記可以代表動量、位置、自旋，或任何完備的單粒子基底。由於產生算符彼此對易，
\begin{equation}
A_i^\dagger A_j^\dagger\ket{0}
=A_j^\dagger A_i^\dagger\ket{0}.
\end{equation}
雙玻色子態在交換兩個粒子時保持對稱。

費米子算符滿足反對易關係：
\begin{align}
\{C_i,C_j\}&=0,&
\{C_i^\dagger,C_j^\dagger\}&=0,&
\{C_i,C_j^\dagger\}&=\delta_{ij},
\end{align}
其中
\begin{equation}
\{X,Y\}=XY+YX.
\end{equation}
對單一模式而言，
\begin{equation}
\{C^\dagger,C^\dagger\}=2(C^\dagger)^2=0,
\end{equation}
因此
\begin{equation}
(C^\dagger)^2=0.
\end{equation}
這個算符本身不為零，但它的平方為零。它可以在該模式中產生一個費米子；再作用一次就會得到零。這就是泡利不相容原理的算符表述。

對不同的標記而言，
\begin{equation}
C_i^\dagger C_j^\dagger\ket{0}
=-C_j^\dagger C_i^\dagger\ket{0},
\end{equation}
這體現了雙費米子態的反對稱性。代數上的替換關係
\begin{equation}
\begin{gathered}
\text{bosons: commutators}\\
\Updownarrow\\
\text{fermions: anticommutators}
\end{gathered}
\end{equation}
則編碼了統計性質的差異。

許多玻色子可以佔據同一個模式。因此，當光子佔據數很大時，其行為可以類似具有普通數值振幅的經典電磁場。費米子不能在同一模式中不斷累積，而奇費米場也不會取得普通的可交換經典值。因此，描述費米子時需要採用不同的算術。

\section{格拉斯曼數與宇稱}

普通數彼此可交換：
\begin{equation}
\alpha_i\alpha_j=\alpha_j\alpha_i.
\end{equation}
格拉斯曼生成元彼此反交換：
\begin{equation}
\theta_i\theta_j=-\theta_j\theta_i.
\end{equation}
令\(i=j\)，並考慮滿足\(2\neq0\)的數系，便得到
\begin{equation}
\theta_i^2=0.
\end{equation}
\(\theta_i\)不能取六或七這類普通數值。這些符號構成一種用來追蹤費米子符號的代數。

普通係數與格拉斯曼生成元彼此可交換：
\begin{equation}
\alpha\theta_i=\theta_i\alpha.
\end{equation}

\begin{classroomqa}
\audiencequestion{普通數乘上格拉斯曼數會得到什麼類型的物件？\lecturetimestamp{01:12:32}}
\lecturerresponse{結果仍是格拉斯曼奇元素。乘上普通係數只會改變尺度，不會改變宇稱。這些生成元所張成的空間，是以普通數為係數的向量空間。}
\end{classroomqa}

這個代數具有\(\mathbb Z_2\)宇稱。普通係數與偶數個生成元的乘積都是偶元素；單一生成元或奇數個生成元的乘積則是奇元素：
\begin{equation}
\begin{aligned}
\text{even}\cdot\text{even}&=\text{even},\\
\text{even}\cdot\text{odd}&=\text{odd},\\
\text{odd}\cdot\text{odd}&=\text{even}.
\end{aligned}
\end{equation}

\begin{classroomqa}
\audiencequestion{這類似實數與虛數的乘法規則嗎？或者，它像引入\(i\)那樣，是一種推廣嗎？\lecturetimestamp{01:14:25}}
\lecturerresponse{這是算術的一種推廣，但不宜過度延伸這個類比。虛數仍然可交換，並滿足\(i^2=-1\)；格拉斯曼奇量則反交換，且平方為零。兩者雖有相同的奇偶乘法模式，代數結構卻並不相同。}
\end{classroomqa}

像\(\theta_1\theta_2\)這樣的乘積是偶元素，並且與格拉斯曼代數中的每個齊次元素可交換，但它不是普通數。它本身不為零，然而
\begin{equation}
(\theta_1\theta_2)^2=0.
\end{equation}

\begin{classroomqa}
\audiencequestion{既然\(\theta_1\theta_2\)是偶的，為什麼它不只是普通數？\lecturetimestamp{01:17:36}}
\lecturerresponse{偶宇稱表示它在分次代數中可交換，但不表示它是實數或複數純量。它具有冪零性，立刻就能與非零普通數區分開來。}
\end{classroomqa}

\section{格拉斯曼變數的函數}

只有一個生成元時，任何多項式都只會有零階與一階項：
\begin{equation}
f(\theta)=\alpha+\beta\theta.
\end{equation}
所有符合\(n\geq2\)的冪\(\theta^n\)都為零。有兩個生成元時，最一般的元素為
\begin{equation}
f(\theta_1,\theta_2)
=\alpha+\beta_1\theta_1+\beta_2\theta_2
+\gamma\theta_1\theta_2.
\end{equation}
將最後一個單項式的因子順序反轉，符號也會改變：
\begin{equation}
\theta_2\theta_1=-\theta_1\theta_2.
\end{equation}

\begin{classroomqa}
\audiencequestion{將\(\theta_1\)乘以係數\(\beta\)，會產生第三個獨立的格拉斯曼變數嗎？\lecturetimestamp{01:18:39}}
\lecturerresponse{不會。這只會得到\(\theta_1\)所張成空間中的另一個向量，並不會產生新的基底生成元。只有兩個獨立生成元時，所有奇線性項都可寫成\(\beta_1\theta_1+\beta_2\theta_2\)。}
\end{classroomqa}

三個生成元的有序基底為
\begin{equation}
1,\quad
\theta_1,\theta_2,\theta_3,\quad
\theta_1\theta_2,\theta_1\theta_3,\theta_2\theta_3,\quad
\theta_1\theta_2\theta_3.
\end{equation}
一般而言，\(n\)個生成元會產生\(2^n\)個有序基底單項式。有限個格拉斯曼生成元的任何解析函數都會截斷成這類多項式，因為生成元重複出現時乘積為零。\editorialnote{\(2^n\)維基底明確呈現了講座以一、二、三個生成元建構出的規律。}

\section{格拉斯曼微分}

對單一變數，定義
\begin{equation}
\frac{\partial}{\partial\theta}\alpha=0,
\qquad
\frac{\partial}{\partial\theta}\theta=1,
\end{equation}
使得
\begin{equation}
\frac{\partial}{\partial\theta}
(\alpha+\beta\theta)=\beta.
\end{equation}
對兩個變數，採用左導數，並保持標準順序\(\theta_1\theta_2\)：
\begin{align}
\frac{\partial^L f}{\partial\theta_1}
&=\beta_1+\gamma\theta_2,\\
\frac{\partial^L f}{\partial\theta_2}
&=\beta_2-\gamma\theta_1.
\end{align}
第二行的負號來自奇導數必須先越過\(\theta_1\)，才能對\(\theta_2\)微分。

這可由分次乘法法則表達。若\(A\)的宇稱為\(|A|\in\{0,1\}\)，則
\begin{equation}
\frac{\partial^L}{\partial\theta}(AB)
=\left(\frac{\partial^L A}{\partial\theta}\right)B
+(-1)^{|A|}A
\left(\frac{\partial^L B}{\partial\theta}\right).
\end{equation}
這個法則的形式大致與普通乘法法則相同，但多了反交換性所要求的符號。格拉斯曼積分也可以定義；由於函數空間是有限維的，它可歸結為一個簡短的規則表。講座將那個表留待之後介紹。

\section{指數謎題的解答}

若只有一個格拉斯曼生成元，且普通係數為\(a\)，則
\begin{equation}
e^{a\theta}
=1+a\theta,
\end{equation}
因為所有更高次冪都為零。含有同一生成元的兩個指數滿足
\begin{align}
e^{a\theta}e^{b\theta}
&=(1+a\theta)(1+b\theta)\\
&=1+(a+b)\theta\\
&=e^{(a+b)\theta}.
\end{align}

接著，大家嘗試使用不同的生成元。直接相乘可得
\begin{equation}
e^{\theta_1}e^{\theta_2}
=(1+\theta_1)(1+\theta_2)
=1+\theta_1+\theta_2+\theta_1\theta_2.
\end{equation}
一開始，黑板上的展開式似乎暗示\(\theta_1+\theta_2\)的平方會給出最後一項。此時一位聽眾指出了符號錯誤：
\begin{equation}
(\theta_1+\theta_2)^2
=\theta_1^2+\theta_1\theta_2
 +\theta_2\theta_1+\theta_2^2
=0.
\end{equation}

\begin{classroomqa}
\audiencequestion{平方項會不會因為含有\(\theta_1\theta_2+\theta_2\theta_1\)而抵消？\lecturetimestamp{01:24:00}}
\lecturerresponse{會。這項修正改變了結論：\(e^{\theta_1+\theta_2}=1+\theta_1+\theta_2\)，而有序乘積\(e^{\theta_1}e^{\theta_2}\)也含有\(\theta_1\theta_2\)。樸素的指數加法公式不適用於不同的反交換指數。}
\end{classroomqa}

這並不矛盾。即使是普通矩陣，也只有在\(A\)與\(B\)對易時，才滿足\(e^Ae^B=e^{A+B}\)。此處
\begin{equation}
[\theta_1,\theta_2]
=\theta_1\theta_2-\theta_2\theta_1
=2\theta_1\theta_2.
\end{equation}
這個交換子為偶，且與兩個生成元都對易，因此 Baker–Campbell–Hausdorff 級數會終止：
\begin{equation}
e^{\theta_1}e^{\theta_2}
=\exp\left(
\theta_1+\theta_2+\frac12[\theta_1,\theta_2]
\right)
=e^{\theta_1+\theta_2+\theta_1\theta_2}.
\end{equation}
右側的指數是冪零的，因此恰好展開為
\(1+\theta_1+\theta_2+\theta_1\theta_2\)。\editorialnote{這項 BCH 計算解答了講座現場留下的問題。講座正確地發現樸素公式不成立，但在指出缺少的交換子項之前便已結束。}

\begin{classroomqa}
\audiencequestion{\(1+\theta_1\)是什麼類型的數？普通部分和格拉斯曼部分可以相加嗎？\lecturetimestamp{01:26:16}}
\lecturerresponse{這是格拉斯曼代數中的非齊次元素，包含偶部分\(1\)與奇部分\(\theta_1\)。這類和是允許的，但既不是普通純量，也不是齊次奇元素。}
\end{classroomqa}

接近結尾時，講座又檢驗了另一個看似矛盾的結果：
\begin{equation}
(\theta_1+\theta_2)^2=0.
\end{equation}
這不代表\(\theta_1\theta_2=0\)，因為兩個交叉項是彼此抵消，而不是各自為零。

\begin{classroomqa}
\audiencequestion{如果每個奇 Grassmann 元素的平方都為零，是否表示任意兩個 Grassmann 生成元的乘積也都為零？\lecturetimestamp{01:32:19}}
\lecturerresponse{不是。方程式 \((\theta_1+\theta_2)^2=0\) 表示 \(\theta_1\theta_2+\theta_2\theta_1=0\)。這證明的是反對易性，並不表示任一種排列下的乘積為零。不同的生成元相乘，結果可以不為零。}
\end{classroomqa}

\begin{classroomqa}
\audiencequestion{我們寫下這些指數函數時，是否假設採用通常的冪級數定義？\lecturetimestamp{01:34:42}}
\lecturerresponse{是的，而且這個定義是自洽的，因為每個級數都會終止。失效的不是指數函數本身，而是對非對易的 \(A\) 與 \(B\)，未經論證便套用恆等式 \(e^Ae^B=e^{A+B}\)。對易子修正項使合成律恢復正確。}
\end{classroomqa}

這樣以探索性的方式收尾恰如其分。從普通對稱性出發，我們由酉變換走到生成元與對易子；從超對稱性出發，我們由奇生成元走到反對易變數。在這兩種情形中，乘法律都是核心。對轉動而言，它體現在 \([L_i,L_j]\)；對 Grassmann 指數而言，它則體現在阻止直接相加的對易子項中。這套代數雖然陌生，卻自洽而有限，恰好適用於超對稱性所要求的費米子變換。
